\documentclass[11pt,a4paper]{article}
\usepackage[utf8]{inputenc}
\usepackage[T1]{fontenc}
\usepackage[catalan]{babel}
\usepackage{lmodern}
\usepackage[sfdefault,lining]{FiraSans}
\usepackage[a4paper,top=2.3cm,bottom=2.2cm,left=1.8cm,right=1.8cm,headheight=1.3cm,headsep=0.4cm,footskip=1cm]{geometry}
\usepackage{amsmath,amssymb}
\usepackage{icomma}
\usepackage{graphicx}
\usepackage[table]{xcolor}
\usepackage{tikz}
\usepackage{fontawesome5}
\usepackage{ccicons}
\usepackage{fancyhdr}
\usepackage{lastpage}
\usepackage{tasks}
\usepackage{needspace}
\usepackage[most]{tcolorbox}
\usepackage{array}
\usepackage[hidelinks,pdftitle={Fitxa d'exercicis de la unitat 8: Àrees i perímetres. 1r d'ESO},pdfsubject={Matemàtiques 1r ESO},pdfkeywords={perímetre, àrea, unitats de longitud, unitats de superfície, hectàrea, quadrat, rectangle, romboide, triangle, rombe, trapezi, polígons regulars, apotema, circumferència, nombre pi, cercle, sector circular, corona circular, figures compostes}]{hyperref}

\definecolor{navy}{HTML}{1B2748}
\definecolor{accent}{HTML}{E8603A}
\definecolor{teal}{HTML}{1F9E90}
\definecolor{soft}{HTML}{EEF1F8}
\definecolor{muted}{HTML}{5D6685}

\setlength{\parindent}{0pt}
\setlength{\parskip}{3pt}
\settasks{label=\alph*), label-format=\bfseries\color{navy}, label-width=1.3em, item-indent=1.8em, column-sep=0.6em, after-item-skip=5pt, before-skip=3pt, after-skip=4pt}

% ---------- capçalera i peu ----------
\newcommand{\llicencia}{\href{https://creativecommons.org/licenses/by-nc-sa/4.0/deed.ca}{\ccbyncsa}\ \ Llicència Creative Commons Reconeixement-NoComercial-CompartirIgual 4.0 Internacional (CC BY-NC-SA 4.0)}
\pagestyle{fancy}
\fancyhf{}
\fancyhead[L]{\raisebox{-0.35\height}{\includegraphics[height=0.95cm]{logo-gva.png}}}
\fancyhead[C]{\small\color{muted}Fitxa d'exercicis · Unitat 8 · 1r d'ESO}
\fancyhead[R]{\raisebox{-0.35\height}{\includegraphics[height=0.95cm]{logo-vermellar.png}}}
\fancyfoot[L]{\scriptsize\color{muted}\llicencia}
\fancyfoot[R]{\small\color{muted}\thepage}
\renewcommand{\headrulewidth}{0.4pt}
\renewcommand{\headrule}{\color{muted!40}\hrule height \headrulewidth}
\fancypagestyle{primera}{\fancyhf{}\renewcommand{\headrulewidth}{0pt}
  \fancyfoot[L]{\scriptsize\color{muted}\llicencia}
  \fancyfoot[R]{\small\color{muted}\thepage}}

% ---------- nivells i exercicis ----------
\newcommand{\estrelles}[1]{{\color{accent}\ifcase#1\or\faStar\faStar[regular]\faStar[regular]\or\faStar\faStar\faStar[regular]\or\faStar\faStar\faStar\fi}}
\newcounter{exer}
\newcommand{\ex}[1]{\par\needspace{4\baselineskip}\medskip\stepcounter{exer}\noindent{\color{navy}\bfseries\theexer.}\ \ {\small\estrelles{#1}}\ \ }

\newcommand{\seccio}[2]{\par\needspace{10\baselineskip}\bigskip
  \begin{tcolorbox}[colback=navy,colframe=navy,coltext=white,arc=2mm,boxrule=0pt,left=3mm,right=3mm,top=1.2mm,bottom=1.2mm]
  {\large\bfseries #1\quad #2}
  \end{tcolorbox}\vspace{-2pt}}

\newcolumntype{C}[1]{>{\centering\arraybackslash}p{#1}}

% quadre de recordatori (només a la fitxa d'exercicis)
\newcommand{\recorda}[1]{\par\smallskip\begin{tcolorbox}[colback=accent!7,colframe=accent!55,arc=2mm,boxrule=0.5pt,left=2.5mm,right=2.5mm,top=1mm,bottom=1mm]\small{\color{accent}\faLightbulb[regular]}\ \textbf{Recorda:}\ #1\end{tcolorbox}}
\newcommand{\res}[1]{#1}\newcommand{\rest}[1]{#1}

\def\fmneed{14}
% =====================================================================
% Capçalera FentMates de les fitxes: un color per a cada curs.
% Els colors són els de la pàgina inicial (LEVELS a assets/js/materials.js);
% el logo ve de _eines/logo/pdf/ (python3 _eines/logo/genera.py).
%
% Ús, al final del preàmbul de la unitat:
%   % =====================================================================
% Capçalera FentMates de les fitxes: un color per a cada curs.
% Els colors són els de la pàgina inicial (LEVELS a assets/js/materials.js);
% el logo ve de _eines/logo/pdf/ (python3 _eines/logo/genera.py).
%
% Ús, al final del preàmbul de la unitat:
%   % =====================================================================
% Capçalera FentMates de les fitxes: un color per a cada curs.
% Els colors són els de la pàgina inicial (LEVELS a assets/js/materials.js);
% el logo ve de _eines/logo/pdf/ (python3 _eines/logo/genera.py).
%
% Ús, al final del preàmbul de la unitat:
%   \input{../comu/fentmates.tex}
%   \fmcurs{1eso}
% i, al començament del document:
%   \capcalera{Matemàtiques}{1r ESO}{Fitxa d'exercicis · Unitat 2}{Divisibilitat}
% Opcional abans de \input: \def\fmneed{14} (espai mínim abans de cada secció).
% =====================================================================
\makeatletter
\def\fm@logodir{../../logo/pdf}
\providecommand{\fmneed}{10}
\def\fm@cursnom{}

% clau del curs, color i color del text damunt del color
\newcommand{\fm@curs}[3]{\@namedef{fm@color@#1}{#2}\@namedef{fm@text@#1}{#3}}
\fm@curs{1eso}{E6B84A}{navy}
\fm@curs{2eso}{3B82F6}{white}
\fm@curs{3eso}{FF7A50}{white}
\fm@curs{4eso-a}{A78BFA}{white}
\fm@curs{4eso-b}{6D28D9}{white}
\fm@curs{1batx-mi}{9B1C1C}{white}
\fm@curs{1batx-ccss}{EF6B6B}{white}
\fm@curs{2batx-mii}{1E6B47}{white}
\fm@curs{2batx-ccss}{2BB3A3}{white}

\newcommand{\fmcurs}[1]{%
  \@ifundefined{fm@color@#1}{\PackageError{fentmates}{Curs desconegut: #1}{}}{}%
  \edef\fm@tmp{\noexpand\definecolor{curs}{HTML}{\@nameuse{fm@color@#1}}%
               \noexpand\colorlet{curstext}{\@nameuse{fm@text@#1}}}\fm@tmp
  \colorlet{cursfosc}{curs!72!black}%
  \def\fm@logo{\fm@logodir/vertical-#1.pdf}%
  \def\fm@marca{\fm@logodir/marca-#1.pdf}}

% ---------- capçalera de la primera pàgina ----------
% #1 matèria (petit), #2 curs, #3 tipus de document i unitat, #4 títol
\newcommand{\capcalera}[4]{%
  \gdef\fm@cursnom{#2}%
  \par\noindent
  \begin{tcolorbox}[colback=curs,colframe=curs,boxrule=0pt,arc=2.2mm,width=3.5cm,height=2.15cm,
      valign=center,left=3mm,right=2mm,top=1mm,bottom=1mm,nobeforeafter,box align=center]
    \color{curstext}{\footnotesize #1}\par\vspace{1pt}{\fontsize{19}{21}\selectfont\bfseries #2\par}%
  \end{tcolorbox}\hspace{5mm}%
  \begin{minipage}[c]{\dimexpr\linewidth-3.5cm-5mm-2.9cm\relax}
    {\color{cursfosc}\small\bfseries\MakeUppercase{#3}}\par\vspace{2pt}
    {\color{navy}\fontsize{20}{23}\selectfont\bfseries #4\par}
  \end{minipage}\hfill
  \begin{minipage}[c]{2.6cm}\raggedleft\includegraphics[width=2.6cm]{\fm@logo}\end{minipage}\par
  \vspace{2pt}}

% ---------- barres de secció del color del curs ----------
\renewcommand{\seccio}[2]{\par\needspace{\fmneed\baselineskip}\bigskip
  \begin{tcolorbox}[colback=curs,colframe=curs,coltext=curstext,arc=2mm,boxrule=0pt,left=3mm,right=3mm,top=1.2mm,bottom=1.2mm]
  {\large\bfseries #1\quad #2}
  \end{tcolorbox}\vspace{-2pt}}

% ---------- estrelles de dificultat del color del curs (a Batxillerat no se'n posen) ----------
\@ifundefined{estrelles}{}{%
  \renewcommand{\estrelles}[1]{{\color{curs}\ifcase#1\or\faStar\faStar[regular]\faStar[regular]\or\faStar\faStar\faStar[regular]\or\faStar\faStar\faStar\fi}}}

% ---------- capçalera i peu de les pàgines ----------
% primera pàgina: logos institucionals al peu
\fancypagestyle{primera}{\fancyhf{}\renewcommand{\headrulewidth}{0pt}
  \fancyfoot[L]{\raisebox{-0.32\height}{\includegraphics[height=0.78cm]{logo-gva.png}}\hspace{5mm}%
    \raisebox{-0.32\height}{\includegraphics[height=0.78cm]{logo-vermellar.png}}}
  \fancyfoot[R]{\scriptsize\color{muted}\href{https://creativecommons.org/licenses/by-nc-sa/4.0/deed.ca}{\ccbyncsa}\ \ CC BY-NC-SA 4.0 · fentmates.org\qquad\small\thepage}}
% la resta: el símbol a l'esquerra i el curs a la dreta, amb el títol al centre com abans
\fancyhead[L]{\raisebox{-0.3\height}{\includegraphics[height=0.62cm]{\fm@marca}}}
\fancyhead[R]{\small\bfseries\color{cursfosc}\fm@cursnom}
\renewcommand{\headrule}{\color{curs!45}\hrule height \headrulewidth}
\makeatother

%   \fmcurs{1eso}
% i, al començament del document:
%   \capcalera{Matemàtiques}{1r ESO}{Fitxa d'exercicis · Unitat 2}{Divisibilitat}
% Opcional abans de \input: \def\fmneed{14} (espai mínim abans de cada secció).
% =====================================================================
\makeatletter
\def\fm@logodir{../../logo/pdf}
\providecommand{\fmneed}{10}
\def\fm@cursnom{}

% clau del curs, color i color del text damunt del color
\newcommand{\fm@curs}[3]{\@namedef{fm@color@#1}{#2}\@namedef{fm@text@#1}{#3}}
\fm@curs{1eso}{E6B84A}{navy}
\fm@curs{2eso}{3B82F6}{white}
\fm@curs{3eso}{FF7A50}{white}
\fm@curs{4eso-a}{A78BFA}{white}
\fm@curs{4eso-b}{6D28D9}{white}
\fm@curs{1batx-mi}{9B1C1C}{white}
\fm@curs{1batx-ccss}{EF6B6B}{white}
\fm@curs{2batx-mii}{1E6B47}{white}
\fm@curs{2batx-ccss}{2BB3A3}{white}

\newcommand{\fmcurs}[1]{%
  \@ifundefined{fm@color@#1}{\PackageError{fentmates}{Curs desconegut: #1}{}}{}%
  \edef\fm@tmp{\noexpand\definecolor{curs}{HTML}{\@nameuse{fm@color@#1}}%
               \noexpand\colorlet{curstext}{\@nameuse{fm@text@#1}}}\fm@tmp
  \colorlet{cursfosc}{curs!72!black}%
  \def\fm@logo{\fm@logodir/vertical-#1.pdf}%
  \def\fm@marca{\fm@logodir/marca-#1.pdf}}

% ---------- capçalera de la primera pàgina ----------
% #1 matèria (petit), #2 curs, #3 tipus de document i unitat, #4 títol
\newcommand{\capcalera}[4]{%
  \gdef\fm@cursnom{#2}%
  \par\noindent
  \begin{tcolorbox}[colback=curs,colframe=curs,boxrule=0pt,arc=2.2mm,width=3.5cm,height=2.15cm,
      valign=center,left=3mm,right=2mm,top=1mm,bottom=1mm,nobeforeafter,box align=center]
    \color{curstext}{\footnotesize #1}\par\vspace{1pt}{\fontsize{19}{21}\selectfont\bfseries #2\par}%
  \end{tcolorbox}\hspace{5mm}%
  \begin{minipage}[c]{\dimexpr\linewidth-3.5cm-5mm-2.9cm\relax}
    {\color{cursfosc}\small\bfseries\MakeUppercase{#3}}\par\vspace{2pt}
    {\color{navy}\fontsize{20}{23}\selectfont\bfseries #4\par}
  \end{minipage}\hfill
  \begin{minipage}[c]{2.6cm}\raggedleft\includegraphics[width=2.6cm]{\fm@logo}\end{minipage}\par
  \vspace{2pt}}

% ---------- barres de secció del color del curs ----------
\renewcommand{\seccio}[2]{\par\needspace{\fmneed\baselineskip}\bigskip
  \begin{tcolorbox}[colback=curs,colframe=curs,coltext=curstext,arc=2mm,boxrule=0pt,left=3mm,right=3mm,top=1.2mm,bottom=1.2mm]
  {\large\bfseries #1\quad #2}
  \end{tcolorbox}\vspace{-2pt}}

% ---------- estrelles de dificultat del color del curs (a Batxillerat no se'n posen) ----------
\@ifundefined{estrelles}{}{%
  \renewcommand{\estrelles}[1]{{\color{curs}\ifcase#1\or\faStar\faStar[regular]\faStar[regular]\or\faStar\faStar\faStar[regular]\or\faStar\faStar\faStar\fi}}}

% ---------- capçalera i peu de les pàgines ----------
% primera pàgina: logos institucionals al peu
\fancypagestyle{primera}{\fancyhf{}\renewcommand{\headrulewidth}{0pt}
  \fancyfoot[L]{\raisebox{-0.32\height}{\includegraphics[height=0.78cm]{logo-gva.png}}\hspace{5mm}%
    \raisebox{-0.32\height}{\includegraphics[height=0.78cm]{logo-vermellar.png}}}
  \fancyfoot[R]{\scriptsize\color{muted}\href{https://creativecommons.org/licenses/by-nc-sa/4.0/deed.ca}{\ccbyncsa}\ \ CC BY-NC-SA 4.0 · fentmates.org\qquad\small\thepage}}
% la resta: el símbol a l'esquerra i el curs a la dreta, amb el títol al centre com abans
\fancyhead[L]{\raisebox{-0.3\height}{\includegraphics[height=0.62cm]{\fm@marca}}}
\fancyhead[R]{\small\bfseries\color{cursfosc}\fm@cursnom}
\renewcommand{\headrule}{\color{curs!45}\hrule height \headrulewidth}
\makeatother

%   \fmcurs{1eso}
% i, al començament del document:
%   \capcalera{Matemàtiques}{1r ESO}{Fitxa d'exercicis · Unitat 2}{Divisibilitat}
% Opcional abans de \input: \def\fmneed{14} (espai mínim abans de cada secció).
% =====================================================================
\makeatletter
\def\fm@logodir{../../logo/pdf}
\providecommand{\fmneed}{10}
\def\fm@cursnom{}

% clau del curs, color i color del text damunt del color
\newcommand{\fm@curs}[3]{\@namedef{fm@color@#1}{#2}\@namedef{fm@text@#1}{#3}}
\fm@curs{1eso}{E6B84A}{navy}
\fm@curs{2eso}{3B82F6}{white}
\fm@curs{3eso}{FF7A50}{white}
\fm@curs{4eso-a}{A78BFA}{white}
\fm@curs{4eso-b}{6D28D9}{white}
\fm@curs{1batx-mi}{9B1C1C}{white}
\fm@curs{1batx-ccss}{EF6B6B}{white}
\fm@curs{2batx-mii}{1E6B47}{white}
\fm@curs{2batx-ccss}{2BB3A3}{white}

\newcommand{\fmcurs}[1]{%
  \@ifundefined{fm@color@#1}{\PackageError{fentmates}{Curs desconegut: #1}{}}{}%
  \edef\fm@tmp{\noexpand\definecolor{curs}{HTML}{\@nameuse{fm@color@#1}}%
               \noexpand\colorlet{curstext}{\@nameuse{fm@text@#1}}}\fm@tmp
  \colorlet{cursfosc}{curs!72!black}%
  \def\fm@logo{\fm@logodir/vertical-#1.pdf}%
  \def\fm@marca{\fm@logodir/marca-#1.pdf}}

% ---------- capçalera de la primera pàgina ----------
% #1 matèria (petit), #2 curs, #3 tipus de document i unitat, #4 títol
\newcommand{\capcalera}[4]{%
  \gdef\fm@cursnom{#2}%
  \par\noindent
  \begin{tcolorbox}[colback=curs,colframe=curs,boxrule=0pt,arc=2.2mm,width=3.5cm,height=2.15cm,
      valign=center,left=3mm,right=2mm,top=1mm,bottom=1mm,nobeforeafter,box align=center]
    \color{curstext}{\footnotesize #1}\par\vspace{1pt}{\fontsize{19}{21}\selectfont\bfseries #2\par}%
  \end{tcolorbox}\hspace{5mm}%
  \begin{minipage}[c]{\dimexpr\linewidth-3.5cm-5mm-2.9cm\relax}
    {\color{cursfosc}\small\bfseries\MakeUppercase{#3}}\par\vspace{2pt}
    {\color{navy}\fontsize{20}{23}\selectfont\bfseries #4\par}
  \end{minipage}\hfill
  \begin{minipage}[c]{2.6cm}\raggedleft\includegraphics[width=2.6cm]{\fm@logo}\end{minipage}\par
  \vspace{2pt}}

% ---------- barres de secció del color del curs ----------
\renewcommand{\seccio}[2]{\par\needspace{\fmneed\baselineskip}\bigskip
  \begin{tcolorbox}[colback=curs,colframe=curs,coltext=curstext,arc=2mm,boxrule=0pt,left=3mm,right=3mm,top=1.2mm,bottom=1.2mm]
  {\large\bfseries #1\quad #2}
  \end{tcolorbox}\vspace{-2pt}}

% ---------- estrelles de dificultat del color del curs (a Batxillerat no se'n posen) ----------
\@ifundefined{estrelles}{}{%
  \renewcommand{\estrelles}[1]{{\color{curs}\ifcase#1\or\faStar\faStar[regular]\faStar[regular]\or\faStar\faStar\faStar[regular]\or\faStar\faStar\faStar\fi}}}

% ---------- capçalera i peu de les pàgines ----------
% primera pàgina: logos institucionals al peu
\fancypagestyle{primera}{\fancyhf{}\renewcommand{\headrulewidth}{0pt}
  \fancyfoot[L]{\raisebox{-0.32\height}{\includegraphics[height=0.78cm]{logo-gva.png}}\hspace{5mm}%
    \raisebox{-0.32\height}{\includegraphics[height=0.78cm]{logo-vermellar.png}}}
  \fancyfoot[R]{\scriptsize\color{muted}\href{https://creativecommons.org/licenses/by-nc-sa/4.0/deed.ca}{\ccbyncsa}\ \ CC BY-NC-SA 4.0 · fentmates.org\qquad\small\thepage}}
% la resta: el símbol a l'esquerra i el curs a la dreta, amb el títol al centre com abans
\fancyhead[L]{\raisebox{-0.3\height}{\includegraphics[height=0.62cm]{\fm@marca}}}
\fancyhead[R]{\small\bfseries\color{cursfosc}\fm@cursnom}
\renewcommand{\headrule}{\color{curs!45}\hrule height \headrulewidth}
\makeatother

\fmcurs{1eso}

\begin{document}
\thispagestyle{primera}

\capcalera{Matemàtiques}{1r ESO}{Fitxa d'exercicis · Unitat 8}{Àrees i perímetres}
\vspace{10pt}
Nom i cognoms: \rule{8.4cm}{0.4pt}\hfill Grup: \rule{1.6cm}{0.4pt}\hfill Data: \rule{2.2cm}{0.4pt}
\vspace{8pt}

\begin{tcolorbox}[colback=soft,colframe=soft,arc=2mm,left=3mm,right=3mm,top=1.5mm,bottom=1.5mm]
\small Exercicis graduats per nivells de dificultat:\quad \estrelles{1}\ bàsic \qquad \estrelles{2}\ mitjà \qquad \estrelles{3}\ avançat.\\[2pt] Resol-los al quadern mostrant tots els passos i escriu sempre les unitats (cm, cm$^2$…). Quan hi haja cercles o circumferències, fes servir $\pi\approx3,14$.
\end{tcolorbox}

% =====================================================
\seccio{1}{Perímetre i àrea. Unitats de mesura}

\ex{1} Cada quadret fa 1\,cm de costat. Calcula l'àrea i el perímetre de cada figura:\begin{tasks}(3)
  \task \begin{tikzpicture}[x=0.55cm,y=0.55cm,baseline=(current bounding box.center),line join=round,line cap=round]
\fill[teal!25] (0,0) rectangle (1,1);
\fill[teal!25] (1,0) rectangle (2,1);
\fill[teal!25] (2,0) rectangle (3,1);
\fill[teal!25] (0,1) rectangle (1,2);
\fill[teal!25] (1,1) rectangle (2,2);
\fill[teal!25] (2,1) rectangle (3,2);
\draw[muted!80,line width=0.5pt,step=1] (0.000,0.000) grid (3.000,2.000);
\draw[thick] (0.000,0.000) -- (1.000,0.000);
\draw[thick] (0.000,0.000) -- (0.000,1.000);
\draw[thick] (1.000,0.000) -- (2.000,0.000);
\draw[thick] (2.000,0.000) -- (3.000,0.000);
\draw[thick] (3.000,0.000) -- (3.000,1.000);
\draw[thick] (0.000,2.000) -- (1.000,2.000);
\draw[thick] (0.000,1.000) -- (0.000,2.000);
\draw[thick] (1.000,2.000) -- (2.000,2.000);
\draw[thick] (2.000,2.000) -- (3.000,2.000);
\draw[thick] (3.000,1.000) -- (3.000,2.000);
\end{tikzpicture}
  \task \begin{tikzpicture}[x=0.55cm,y=0.55cm,baseline=(current bounding box.center),line join=round,line cap=round]
\fill[teal!25] (0,0) rectangle (1,1);
\fill[teal!25] (1,0) rectangle (2,1);
\fill[teal!25] (2,0) rectangle (3,1);
\fill[teal!25] (3,0) rectangle (4,1);
\fill[teal!25] (3,1) rectangle (4,2);
\fill[teal!25] (3,2) rectangle (4,3);
\fill[teal!25] (0,1) rectangle (1,2);
\draw[muted!80,line width=0.5pt,step=1] (0.000,0.000) grid (4.000,3.000);
\draw[thick] (0.000,0.000) -- (1.000,0.000);
\draw[thick] (0.000,0.000) -- (0.000,1.000);
\draw[thick] (1.000,0.000) -- (2.000,0.000);
\draw[thick] (1.000,1.000) -- (2.000,1.000);
\draw[thick] (2.000,0.000) -- (3.000,0.000);
\draw[thick] (2.000,1.000) -- (3.000,1.000);
\draw[thick] (3.000,0.000) -- (4.000,0.000);
\draw[thick] (4.000,0.000) -- (4.000,1.000);
\draw[thick] (3.000,1.000) -- (3.000,2.000);
\draw[thick] (4.000,1.000) -- (4.000,2.000);
\draw[thick] (3.000,3.000) -- (4.000,3.000);
\draw[thick] (3.000,2.000) -- (3.000,3.000);
\draw[thick] (4.000,2.000) -- (4.000,3.000);
\draw[thick] (0.000,2.000) -- (1.000,2.000);
\draw[thick] (0.000,1.000) -- (0.000,2.000);
\draw[thick] (1.000,1.000) -- (1.000,2.000);
\end{tikzpicture}
  \task \begin{tikzpicture}[x=0.55cm,y=0.55cm,baseline=(current bounding box.center),line join=round,line cap=round]
\fill[teal!25] (1,0) rectangle (2,1);
\fill[teal!25] (0,1) rectangle (1,2);
\fill[teal!25] (1,1) rectangle (2,2);
\fill[teal!25] (2,1) rectangle (3,2);
\fill[teal!25] (1,2) rectangle (2,3);
\draw[muted!80,line width=0.5pt,step=1] (0.000,0.000) grid (3.000,3.000);
\draw[thick] (1.000,0.000) -- (2.000,0.000);
\draw[thick] (1.000,0.000) -- (1.000,1.000);
\draw[thick] (2.000,0.000) -- (2.000,1.000);
\draw[thick] (0.000,1.000) -- (1.000,1.000);
\draw[thick] (0.000,2.000) -- (1.000,2.000);
\draw[thick] (0.000,1.000) -- (0.000,2.000);
\draw[thick] (2.000,1.000) -- (3.000,1.000);
\draw[thick] (2.000,2.000) -- (3.000,2.000);
\draw[thick] (3.000,1.000) -- (3.000,2.000);
\draw[thick] (1.000,3.000) -- (2.000,3.000);
\draw[thick] (1.000,2.000) -- (1.000,3.000);
\draw[thick] (2.000,2.000) -- (2.000,3.000);
\end{tikzpicture}
\end{tasks}\recorda{El \textbf{perímetre} és la longitud del contorn (cm, m…) i l'\textbf{àrea} és la superfície que ocupa la figura (cm$^2$, m$^2$…). En una quadrícula, l'àrea és el nombre de quadrets i el perímetre, el nombre de costats de quadret del contorn.}

\ex{1} Expressa en metres:\begin{tasks}(5)
  \task 4,7\,km
  \task 380\,cm
  \task 25\,dm
  \task 9\,hm
  \task 1250\,mm
\end{tasks}\recorda{Les unitats de longitud van de 10 en 10: km, hm, dam, m, dm, cm, mm. Per a passar a una unitat més menuda multipliquem i per a passar a una més gran dividim.}

\ex{1} Fes els canvis d'unitats:\begin{tasks}(4)
  \task 3,2\,m a cm
  \task 45\,mm a cm
  \task 0,8\,km a dam
  \task 760\,dm a hm
\end{tasks}

\ex{1} Fes els canvis d'unitats de superfície:\begin{tasks}(3)
  \task 5\,dm$^2$ a cm$^2$
  \task 2,4\,m$^2$ a dm$^2$
  \task 35\,000\,cm$^2$ a m$^2$
  \task 0,7\,hm$^2$ a m$^2$
  \task 8\,km$^2$ a hm$^2$
\end{tasks}\recorda{Les unitats de superfície van de 100 en 100: \ $1\text{ m}^2=100\text{ dm}^2=10\,000\text{ cm}^2$.}

\ex{2} Fes els canvis amb unitats agràries ($1\text{ ha}=1\text{ hm}^2$, \ $1\text{ a}=1\text{ dam}^2$):\begin{tasks}(4)
  \task 3,5\,ha a m$^2$
  \task 4500\,m$^2$ a ha
  \task 12\,a a m$^2$
  \task 0,25\,km$^2$ a ha
\end{tasks}

\ex{2} Ordena de menor a major: \ $0,05\text{ km}^2$, \ $4\text{ ha}$, \ $520\text{ dam}^2$, \ $48\,000\text{ m}^2$.

\ex{2} Una finca de 2,4\,ha es ven a 18\,€ el metre quadrat. Quant costa?

\ex{3} Volem fer rectangles de 20\,cm de perímetre amb costats que siguen nombres naturals. Escriu-los tots i calcula'n l'àrea. Quin té l'àrea més gran?


% =====================================================
\seccio{2}{Perímetre dels polígons}

\ex{1} Calcula el perímetre de:\begin{tasks}(2)
  \task un triangle de costats 7,5\,cm, 9\,cm i 11\,cm
  \task un quadrat de 8,5\,cm de costat
  \task un rectangle de 15\,m × 9\,m
  \task un rombe de 6,3\,cm de costat
\end{tasks}\recorda{El perímetre d'un polígon és la suma dels seus costats. Quadrat: $P=4\cdot l$. \ Rectangle: $P=2\cdot b+2\cdot h$. \ Polígon regular de $n$ costats: $P=n\cdot l$.}

\ex{1} Calcula el perímetre d'aquests polígons regulars:\begin{tasks}(2)
  \task pentàgon de 7\,cm de costat
  \task hexàgon de 4,5\,cm de costat
  \task octàgon de 3,2\,cm de costat
  \task decàgon de 12\,mm de costat
\end{tasks}

\ex{1} Calcula el perímetre d'aquest polígon:\par\smallskip\begin{center}\begin{tikzpicture}[x=0.45cm,y=0.45cm,baseline=(current bounding box.center),line join=round,line cap=round]
\draw[fill=teal!12,draw=teal,thick] (0.000,0.000) -- (5.000,0.000) -- (5.000,3.000) -- (2.000,5.500) -- (0.000,3.000) -- cycle;
\node[font=\small,below] at (2.500,0.000) {5 cm};
\node[font=\small,right] at (5.000,1.500) {3 cm};
\node[font=\small,above right] at (3.500,4.250) {4 cm};
\node[font=\small,above left] at (1.000,4.250) {3,5 cm};
\node[font=\small,left] at (0.000,1.500) {3 cm};
\end{tikzpicture}\end{center}

\ex{2} Calcula el costat de:\begin{tasks}(1)
  \task un quadrat de 46\,cm de perímetre
  \task un hexàgon regular de 51\,cm de perímetre
  \task un triangle equilàter de 13,5\,m de perímetre
\end{tasks}

\ex{2} Un rectangle té 64\,cm de perímetre i la base fa 20\,cm. Quant fa l'altura?

\ex{2} Volem tancar un camp rectangular de 120\,m × 85\,m amb tres fils de filferro. Quants metres de filferro calen? Quant costaran si el metre val 0,35\,€?

\ex{3} La base d'un rectangle fa el triple que l'altura i el perímetre fa 96\,cm. Quant fan la base i l'altura?

\ex{3} Cadascun dels dos costats iguals d'un triangle isòsceles fa el doble que la base. Si el perímetre és de 40\,cm, quant fa cada costat?


% =====================================================
\seccio{3}{Àrea del quadrat i del rectangle}

\ex{1} Calcula l'àrea i el perímetre de:\begin{tasks}(2)
  \task un quadrat de 12\,cm de costat
  \task un rectangle de 15\,cm × 8\,cm
  \task un rectangle de 7,5\,m × 4\,m
  \task un quadrat de 2,5\,m de costat
\end{tasks}\recorda{Rectangle: $A=b\cdot h$. \ Quadrat: $A=l\cdot l=l^2$. Les dues mesures han d'estar en la mateixa unitat.}

\ex{1} Un rectangle té 84\,cm$^2$ d'àrea i 12\,cm de base. Quant fan l'altura i el perímetre?

\ex{1} Respon:\begin{tasks}(1)
  \task Un quadrat té 121\,m$^2$ d'àrea. Quant fa el costat? I el perímetre?
  \task Un quadrat té 28\,cm de perímetre. Quina és la seua àrea?
\end{tasks}

\ex{2} Una habitació fa 5,4\,m × 3,5\,m. Quant costarà posar-hi parquet si el metre quadrat val 24\,€?

\ex{2} Volem pintar una paret de 6\,m × 2,5\,m que té una finestra d'1,5\,m × 1,2\,m. Un pot de pintura cobreix 5\,m$^2$. Quants pots hem de comprar?

\ex{2} Quantes rajoles quadrades de 50\,cm de costat calen per a enrajolar un pati de 8\,m × 6\,m?

\ex{3} Un quadrat fa 5\,cm de costat. Si dupliquem el costat, què li passa al perímetre? I a l'àrea?

\ex{3} Escriu tots els rectangles de costats naturals que tenen 36\,cm$^2$ d'àrea. Quin té el perímetre més menut?


% =====================================================
\seccio{4}{Àrea del romboide i del triangle}

\ex{1} Calcula l'àrea dels romboides:\begin{tasks}(2)
  \task base 12\,cm i altura 7\,cm
  \task base 8,5\,m i altura 4\,m
\end{tasks}\recorda{Romboide: $A=b\cdot h$. \ Triangle: $A=\dfrac{b\cdot h}{2}$. L'altura sempre és perpendicular a la base.}

\ex{1} Calcula l'àrea dels triangles:\begin{tasks}(3)
  \task base 14\,cm i altura 9\,cm
  \task base 7\,cm i altura 5\,cm
  \task base 6,4\,m i altura 3,5\,m
\end{tasks}

\ex{1} Calcula l'àrea de cada triangle:\begin{tasks}(3)
  \task \begin{tikzpicture}[x=0.42cm,y=0.42cm,baseline=(current bounding box.center),line join=round,line cap=round]
\draw[fill=teal!12,draw=teal,thick] (0.000,0.000) -- (6.000,0.000) -- (2.000,4.000) -- cycle;
\draw[dashed,muted] (2.000,0.000) -- (2.000,4.000);
\draw[thin] (2.524,0.000) -- (2.524,0.524) -- (2.000,0.524);
\node[font=\small,below] at (3.000,0.000) {6 cm};
\node[font=\small,right] at (2.000,2.000) {4 cm};
\end{tikzpicture}
  \task \begin{tikzpicture}[x=0.38cm,y=0.38cm,baseline=(current bounding box.center),line join=round,line cap=round]
\draw[fill=teal!12,draw=teal,thick] (0.000,0.000) -- (8.000,0.000) -- (0.000,5.000) -- cycle;
\draw[thin] (0.579,0.000) -- (0.579,0.579) -- (0.000,0.579);
\node[font=\small,below] at (4.000,0.000) {8 cm};
\node[font=\small,left] at (0.000,2.500) {5 cm};
\end{tikzpicture}
  \task \begin{tikzpicture}[x=0.38cm,y=0.38cm,baseline=(current bounding box.center),line join=round,line cap=round]
\draw[fill=teal!12,draw=teal,thick] (0.000,0.000) -- (4.000,0.000) -- (6.500,5.000) -- cycle;
\draw[dashed,muted] (4.000,0.000) -- (6.500,0.000);
\draw[dashed,muted] (6.500,0.000) -- (6.500,5.000);
\draw[thin] (5.921,0.000) -- (5.921,0.579) -- (6.500,0.579);
\node[font=\small,below] at (2.000,0.000) {4 cm};
\node[font=\small,right] at (6.500,2.500) {5 cm};
\end{tikzpicture}
\end{tasks}

\ex{2} Els catets d'un triangle rectangle fan 9\,cm i 12\,cm, i la hipotenusa, 15\,cm. Calcula'n l'àrea i el perímetre.

\ex{2} Respon:\begin{tasks}(1)
  \task Un romboide té 96\,cm$^2$ d'àrea i 12\,cm de base. Quant fa l'altura?
  \task Un triangle té 45\,m$^2$ d'àrea i 9\,m d'altura. Quant fa la base?
\end{tasks}

\ex{2} Cada quadret de la quadrícula fa 1\,cm de costat. Calcula l'àrea del triangle.\par\smallskip\begin{center}\begin{tikzpicture}[x=0.42cm,y=0.42cm,baseline=(current bounding box.center),line join=round,line cap=round]
\draw[muted!80,line width=0.5pt,step=1] (0.000,0.000) grid (8.000,5.000);
\draw[fill=accent,fill opacity=0.22,draw=accent,thick] (0.000,0.000) -- (8.000,0.000) -- (3.000,5.000) -- cycle;
\end{tikzpicture}\end{center}

\ex{3} Un triangle té els costats de 13\,cm, 14\,cm i 15\,cm, i l'altura sobre el costat de 14\,cm fa 12\,cm. Calcula'n l'àrea i l'altura sobre el costat de 15\,cm.

\ex{3} Dibuixa dues rectes paral·leles separades 6\,cm i, sobre una d'elles, un segment de 10\,cm. Dibuixa tres triangles diferents que tinguen aquest segment de base i el tercer vèrtex sobre l'altra recta. Quina àrea té cadascun? Per què?


% =====================================================
\seccio{5}{Àrea del rombe i del trapezi}

\ex{1} Calcula l'àrea dels rombes:\begin{tasks}(3)
  \task diagonals 16\,cm i 10\,cm
  \task diagonals 9\,m i 6\,m
  \task diagonals 7,5\,cm i 4\,cm
\end{tasks}\recorda{Rombe: $A=\dfrac{D\cdot d}{2}$. \ Trapezi: $A=\dfrac{(B+b)\cdot h}{2}$.}

\ex{1} Calcula l'àrea dels trapezis:\begin{tasks}(2)
  \task bases 12\,cm i 8\,cm, altura 5\,cm
  \task bases 9\,m i 4\,m, altura 6\,m
  \task bases 10,5\,cm i 6,5\,cm, altura 4\,cm
\end{tasks}

\ex{2} Un rombe té 52\,cm de perímetre i les diagonals fan 24\,cm i 10\,cm. Quant fa el costat? Calcula'n l'àrea.

\ex{2} Calcula l'àrea i el perímetre d'aquest trapezi isòsceles:\par\smallskip\begin{center}\begin{tikzpicture}[x=0.3cm,y=0.3cm,baseline=(current bounding box.center),line join=round,line cap=round]
\draw[fill=teal!12,draw=teal,thick] (0.000,0.000) -- (18.000,0.000) -- (14.000,3.000) -- (4.000,3.000) -- cycle;
\draw[dashed,muted] (4.000,0.000) -- (4.000,3.000);
\draw[thin] (4.733,0.000) -- (4.733,0.733) -- (4.000,0.733);
\node[font=\small,below] at (9.000,0.000) {18 cm};
\node[font=\small,above] at (9.000,3.000) {10 cm};
\node[font=\small,right] at (4.000,1.500) {3 cm};
\node[font=\small,right] at (16.000,1.500) {5 cm};
\end{tikzpicture}\end{center}

\ex{2} Respon:\begin{tasks}(1)
  \task Un rombe té 60\,cm$^2$ d'àrea i la diagonal major fa 15\,cm. Quant fa la menor?
  \task Un trapezi té 72\,m$^2$ d'àrea i les bases fan 14\,m i 10\,m. Quant fa l'altura?
\end{tasks}

\ex{3} Un estel té forma de rombe, amb diagonals d'1,2\,m i 80\,cm. Quants m$^2$ de tela calen per a fer-lo? Expressa-ho també en cm$^2$.

\ex{3} Una parcel·la té forma de trapezi, amb bases de 85\,m i 55\,m i 40\,m d'amplària (altura). Quantes hectàrees fa? Quant val si el metre quadrat es paga a 12\,€?


% =====================================================
\seccio{6}{Àrea dels polígons regulars i dels polígons qualssevol}

\ex{1} Un pentàgon regular té 10\,cm de costat i 6,88\,cm d'apotema. Calcula'n el perímetre i l'àrea.\recorda{Polígon regular: $A=\dfrac{P\cdot a}{2}$, on $P=n\cdot l$ és el perímetre i $a$ és l'apotema (el segment del centre al punt mitjà d'un costat).}

\ex{1} Un hexàgon regular té 8\,cm de costat i 6,93\,cm d'apotema. Calcula'n el perímetre i l'àrea.

\ex{2} Un octàgon regular té 40\,cm de perímetre i 6,04\,cm d'apotema. Quant fa el costat? I l'àrea?

\ex{2} Un hexàgon regular té 36\,cm de perímetre i 93,6\,cm$^2$ d'àrea. Quant fa l'apotema?

\ex{2} Cada quadret fa 1\,cm de costat. Calcula l'àrea del polígon descomponent-lo en un rectangle i un triangle.\par\smallskip\begin{center}\begin{tikzpicture}[x=0.4cm,y=0.4cm,baseline=(current bounding box.center),line join=round,line cap=round]
\draw[muted!80,line width=0.5pt,step=1] (0.000,0.000) grid (6.000,6.000);
\draw[fill=accent,fill opacity=0.22,draw=accent,thick] (0.000,0.000) -- (6.000,0.000) -- (6.000,4.000) -- (3.000,6.000) -- (0.000,4.000) -- cycle;
\end{tikzpicture}\end{center}

\ex{3} Cada quadret fa 1\,cm de costat. Calcula l'àrea del triangle completant-lo fins a un rectangle i restant els triangles que sobren.\par\smallskip\begin{center}\begin{tikzpicture}[x=0.4cm,y=0.4cm,baseline=(current bounding box.center),line join=round,line cap=round]
\draw[muted!80,line width=0.5pt,step=1] (0.000,0.000) grid (7.000,6.000);
\draw[fill=accent,fill opacity=0.22,draw=accent,thick] (0.000,0.000) -- (7.000,2.000) -- (3.000,6.000) -- cycle;
\end{tikzpicture}\end{center}

\ex{3} Volem enrajolar 25\,m$^2$ amb rajoles en forma d'hexàgon regular de 10\,cm de costat i 8,66\,cm d'apotema. Quantes rajoles calen, com a mínim?


% =====================================================
\seccio{7}{La circumferència i la seua longitud}

\ex{1} Com s'anomenen els elements $a$, $b$, $c$ i $d$ de la figura?\par\smallskip\begin{center}\begin{tikzpicture}[x=0.75cm,y=0.75cm,baseline=(current bounding box.center),line join=round,line cap=round]
\draw[draw=teal,thick] (0.000,0.000) circle (2);
\draw[thick] (0.000,0.000) -- (-1.286,1.532);
\draw[thick] (-1.879,-0.684) -- (1.879,0.684);
\draw[thick] (-0.684,-1.879) -- (1.532,-1.286);
\draw[accent,very thick] (1.532,1.286) arc[start angle=40,end angle=100,radius=2];
\fill (0.000,0.000) circle (1.5pt);
\node[font=\small,inner sep=1pt] at (0.250,-0.300) {$O$};
\node[font=\small,above right] at (-0.643,0.766) {$a$};
\node[font=\small,below right] at (0.000,0.000) {$b$};
\node[font=\small,above] at (0.424,-1.582) {$c$};
\node[font=\small,inner sep=1pt] at (0.804,2.208) {$d$};
\end{tikzpicture}\end{center}

\ex{1} Calcula la longitud de les circumferències (fes servir $\pi\approx3,14$):\begin{tasks}(4)
  \task radi 4\,cm
  \task radi 9\,m
  \task diàmetre 20\,cm
  \task diàmetre 7\,mm
\end{tasks}\recorda{La longitud de la circumferència és $L=\pi\cdot d=2\cdot\pi\cdot r$, amb $\pi\approx3,14$. Un arc de $n^\circ$ fa $L_{\text{arc}}=\dfrac{2\cdot\pi\cdot r\cdot n}{360}$.}

\ex{2} Respon:\begin{tasks}(1)
  \task Una circumferència fa 31,4\,cm de longitud. Quant fa el radi?
  \task Una circumferència fa 75,36\,m de longitud. Quant fan el diàmetre i el radi?
\end{tasks}

\ex{2} La roda d'una bicicleta fa 70\,cm de diàmetre. Quants centímetres avança en una volta? Quantes voltes ha de fer, aproximadament, per a recórrer 1\,km?

\ex{2} Calcula la longitud dels arcs:\begin{tasks}(3)
  \task $90^\circ$ i radi 8\,cm
  \task $60^\circ$ i radi 12\,cm
  \task $150^\circ$ i radi 6\,cm
\end{tasks}

\ex{3} Calcula el perímetre de:\begin{tasks}(2)
  \task un semicercle de 10\,cm de diàmetre
  \task un quart de cercle de 4\,cm de radi
\end{tasks}

\ex{3} L'agulla dels minuts d'un rellotge fa 15\,cm. Quina distància recorre la punta en 20 minuts? I en una hora?


% =====================================================
\seccio{8}{L'àrea del cercle i de les figures circulars}

\ex{1} Calcula l'àrea dels cercles ($\pi\approx3,14$):\begin{tasks}(4)
  \task radi 6\,cm
  \task radi 2,5\,m
  \task diàmetre 14\,cm
  \task diàmetre 30\,cm
\end{tasks}\recorda{Cercle: $A=\pi\cdot r^2$. \ Sector circular de $n^\circ$: $A=\dfrac{\pi\cdot r^2\cdot n}{360}$. \ Corona circular: $A=\pi\cdot R^2-\pi\cdot r^2$. Si et donen el diàmetre, calcula primer el radi.}

\ex{2} Calcula l'àrea de:\begin{tasks}(2)
  \task un semicercle de 10\,cm de radi
  \task un quart de cercle de 6\,cm de radi
\end{tasks}

\ex{2} Calcula l'àrea dels sectors circulars:\begin{tasks}(3)
  \task $45^\circ$ i radi 8\,cm
  \task $120^\circ$ i radi 9\,cm
  \task $270^\circ$ i radi 4\,cm
\end{tasks}

\ex{2} Calcula l'àrea de la corona circular de la figura en cada cas:\par\smallskip\begin{center}\begin{tikzpicture}[x=0.6cm,y=0.6cm,baseline=(current bounding box.center),line join=round,line cap=round]
\draw[fill=accent!22,draw=accent,thick] (0.000,0.000) circle (2.1);
\draw[fill=white,draw=teal,thick] (0.000,0.000) circle (1.2);
\fill (0.000,0.000) circle (1.5pt);
\draw[thick] (0.000,0.000) -- (1.819,1.050);
\draw[dashed,muted] (0.000,0.000) -- (-1.128,0.410);
\node[font=\small,inner sep=1pt] at (1.226,1.104) {$R$};
\node[font=\small,inner sep=1pt] at (-0.600,0.021) {$r$};
\end{tikzpicture}\end{center}\begin{tasks}(2)
  \task $R=7$\,cm, \ $r=4$\,cm
  \task $R=10$\,m, \ $r=6$\,m
\end{tasks}

\ex{2} Respon:\begin{tasks}(1)
  \task Un cercle té 78,5\,cm$^2$ d'àrea. Quant fa el radi?
  \task Una circumferència fa 18,84\,cm. Quina és l'àrea del cercle que tanca?
\end{tasks}

\ex{3} Un CD fa 12\,cm de diàmetre i té un forat central d'1,5\,cm de diàmetre. Quina és la superfície del CD?

\ex{3} Què té més superfície: una pizza de 32\,cm de diàmetre o dues pizzes de 22\,cm de diàmetre?


% =====================================================
\seccio{9}{Figures compostes}

\ex{1} Calcula l'àrea i el perímetre d'aquesta figura:\par\smallskip\begin{center}\begin{tikzpicture}[x=0.32cm,y=0.32cm,baseline=(current bounding box.center),line join=round,line cap=round]
\draw[fill=teal!12,draw=teal,thick] (0.000,0.000) -- (10.000,0.000) -- (10.000,4.000) -- (6.000,4.000) -- (6.000,7.000) -- (0.000,7.000) -- cycle;
\node[font=\small,below] at (5.000,0.000) {10 cm};
\node[font=\small,left] at (0.000,3.500) {7 cm};
\node[font=\small,right] at (10.000,2.000) {4 cm};
\node[font=\small,above] at (8.000,4.000) {4 cm};
\node[font=\small,above] at (3.000,7.000) {6 cm};
\end{tikzpicture}\end{center}\recorda{Per a l'àrea d'una figura composta, la \textbf{descomponem} en figures conegudes i sumem, o la \textbf{completem} i restem. El perímetre és només el contorn exterior.}

\ex{1} Calcula l'àrea de la façana d'aquesta caseta:\par\smallskip\begin{center}\begin{tikzpicture}[x=0.3cm,y=0.3cm,baseline=(current bounding box.center),line join=round,line cap=round]
\draw[fill=teal!12,draw=teal,thick] (0.000,0.000) -- (8.000,0.000) -- (8.000,5.000) -- (0.000,5.000) -- cycle;
\draw[fill=accent!22,draw=accent,thick] (0.000,5.000) -- (8.000,5.000) -- (4.000,8.000) -- cycle;
\draw[dashed,muted] (4.000,5.000) -- (4.000,8.000);
\node[font=\small,below] at (4.000,0.000) {8 m};
\node[font=\small,right] at (8.000,2.500) {5 m};
\node[font=\small,right] at (4.000,6.500) {3 m};
\end{tikzpicture}\end{center}

\ex{2} Calcula l'àrea i el perímetre de la figura, formada per un rectangle i un semicercle ($\pi\approx3,14$):\par\smallskip\begin{center}\begin{tikzpicture}[x=0.38cm,y=0.38cm,baseline=(current bounding box.center),line join=round,line cap=round]
\draw[fill=accent!22,draw=accent,thick] (0.000,0.000) -- (6.000,0.000) -- (6.000,4.000) -- (0.000,4.000) -- cycle;
\draw[fill=accent!22,draw=accent,thick] (6.000,2.000) -- (6.000,0.000) arc[start angle=-90,end angle=90,radius=2] -- cycle;
\draw[dashed,muted] (6.000,0.000) -- (6.000,4.000);
\node[font=\small,below] at (3.000,0.000) {6 cm};
\node[font=\small,left] at (0.000,2.000) {4 cm};
\end{tikzpicture}\end{center}

\ex{2} Calcula l'àrea de la zona acolorida (un quadrat de 8\,cm de costat sense quatre quarts de cercle de 4\,cm de radi):\par\smallskip\begin{center}\begin{tikzpicture}[x=0.28cm,y=0.28cm,baseline=(current bounding box.center),line join=round,line cap=round]
\draw[fill=accent!22,draw=accent,thick] (0.000,0.000) -- (8.000,0.000) -- (8.000,8.000) -- (0.000,8.000) -- cycle;
\draw[fill=white,draw=teal,thick] (0.000,0.000) -- (4.000,0.000) arc[start angle=0,end angle=90,radius=4] -- cycle;
\draw[fill=white,draw=teal,thick] (8.000,0.000) -- (8.000,4.000) arc[start angle=90,end angle=180,radius=4] -- cycle;
\draw[fill=white,draw=teal,thick] (8.000,8.000) -- (4.000,8.000) arc[start angle=180,end angle=270,radius=4] -- cycle;
\draw[fill=white,draw=teal,thick] (0.000,8.000) -- (-0.000,4.000) arc[start angle=270,end angle=360,radius=4] -- cycle;
\node[font=\small,below] at (4.000,0.000) {8 cm};
\end{tikzpicture}\end{center}

\ex{2} Al voltant d'una piscina de 10\,m × 6\,m es fa una vorera d'1,5\,m d'amplària. Quina és l'àrea de la vorera?\par\smallskip\begin{center}\begin{tikzpicture}[x=0.3cm,y=0.3cm,baseline=(current bounding box.center),line join=round,line cap=round]
\draw[fill=accent!22,draw=accent,thick] (-1.500,-1.500) -- (11.500,-1.500) -- (11.500,7.500) -- (-1.500,7.500) -- cycle;
\draw[fill=white,draw=teal,thick] (0.000,0.000) -- (10.000,0.000) -- (10.000,6.000) -- (0.000,6.000) -- cycle;
\node[font=\small,inner sep=1pt] at (5.000,3.000) {10 m × 6 m};
\node[font=\small,inner sep=1pt] at (5.000,6.750) {1,5 m};
\end{tikzpicture}\end{center}

\ex{2} Una pista està formada per un rectangle de 80\,m × 50\,m i dos semicercles de 50\,m de diàmetre als extrems. Calcula'n l'àrea i la longitud d'una volta.

\ex{3} Una diana té tres circumferències concèntriques de radis 3, 6 i 9\,cm. Calcula l'àrea de cada zona.\par\smallskip\begin{center}\begin{tikzpicture}[x=0.55cm,y=0.55cm,baseline=(current bounding box.center),line join=round,line cap=round]
\draw[fill=accent!22,draw=accent,thick] (0.000,0.000) circle (2.7);
\draw[fill=white,draw=teal,thick] (0.000,0.000) circle (1.8);
\draw[fill=accent!8,draw=accent!60] (0.000,0.000) circle (1.8);
\draw[fill=accent!22,draw=accent,thick] (0.000,0.000) circle (0.9);
\fill (0.000,0.000) circle (1.5pt);
\draw[thick] (0.000,0.000) -- (2.700,0.000);
\node[font=\small,above] at (0.450,0.000) {3};
\node[font=\small,above] at (1.350,0.000) {3};
\node[font=\small,above] at (2.250,0.000) {3};
\end{tikzpicture}\end{center}

\ex{3} Dins d'un quadrat de 6\,cm de costat es dibuixen dos quarts de circumferència amb centre en dos vèrtexs oposats. Calcula l'àrea de la fulla que queda entre els dos arcs.\par\smallskip\begin{center}\begin{tikzpicture}[x=0.38cm,y=0.38cm,baseline=(current bounding box.center),line join=round,line cap=round]
\draw[thick,muted] (0.000,0.000) -- (6.000,0.000) -- (6.000,6.000) -- (0.000,6.000) -- cycle;
\draw[accent,very thick] (6.000,0.000) arc[start angle=0,end angle=90,radius=6];
\draw[accent,very thick] (0.000,6.000) arc[start angle=180,end angle=270,radius=6];
\draw[dashed,muted] (6.000,0.000) -- (0.000,6.000);
\node[font=\small,below] at (3.000,0.000) {6 cm};
\end{tikzpicture}\end{center}

\ex{3} Una finestra està formada per un rectangle d'1,2\,m × 1,5\,m acabat per dalt amb un semicercle. Quant costa el vidre si el m$^2$ val 45\,€?\par\smallskip\begin{center}\begin{tikzpicture}[x=1.6cm,y=1.6cm,baseline=(current bounding box.center),line join=round,line cap=round]
\draw[fill=teal!12,draw=teal,thick] (0.000,0.000) -- (1.200,0.000) -- (1.200,1.500) -- (0.000,1.500) -- cycle;
\draw[fill=teal!12,draw=teal,thick] (0.600,1.500) -- (1.200,1.500) arc[start angle=0,end angle=180,radius=0.6] -- cycle;
\draw[dashed,muted] (0.000,1.500) -- (1.200,1.500);
\node[font=\small,below] at (0.600,0.000) {1,2 m};
\node[font=\small,right] at (1.200,0.750) {1,5 m};
\end{tikzpicture}\end{center}

\vfill
\begin{center}\color{muted}\small\faCheckCircle[regular]\ \ Ànim! Si et costa algun exercici, repassa els apunts interactius de la unitat.\end{center}

\end{document}
